\section{Results}
\subsection{Modern VisDrone}
Table~\ref{tab:modern-vis} reports the aggregate results on the documented
seven-sequence second validation read. For ByteTrack, AC--MOT v3 changes HOTA
from 41.586 to 44.523, MOTA from 9.876 to 27.871, IDF1 from 48.569 to
53.929, and IDS from 1341 to 1016. FP falls from 39274 to 24892, while FN
rises from 24121 to 25902. This is a selective observation stream: clutter and
identity-confusing candidates decrease at a recall cost.

For OATrack, HOTA changes from 46.770 to 48.095, MOTA from 28.165 to 34.483,
IDF1 from 57.107 to 59.364, and IDS from 773 to 604. FP falls from 25394 to
21400 and FN changes from 25432 to 25057. The smaller increment is consistent
with a stronger baseline and a different association response, not evidence
that the controller is ineffective. The two hosts are therefore reported
separately rather than pooled into one AC--MOT gain.

The paired bootstrap in Table~\ref{tab:modern-bootstrap} supports the main
positive differences. ByteTrack HOTA, MOTA, IDF1, and IDS changes are
$+2.937$ [2.356,3.375], $+17.995$ [13.541,23.282], $+5.360$
[4.490,6.178], and $-325$ [-498,-164]. OATrack changes are $+1.325$
[0.200,2.264], $+6.318$ [4.741,7.879], $+2.258$ [0.548,3.474], and
$-169$ [-292,-42]. The ByteTrack FN increase has CI [288,3594.075], whereas
OATrack FN has CI [-1462.575,670]; the evidence is consequently a quality
trade-off rather than an unqualified recall improvement.

The state occupancy is reported in Online Resource~1, Sec.~S2. It confirms
that both hosts use multiple actions and that their trajectories differ, but
occupancy alone does not establish timing causality. The historical V1/V2
records and the pre-v3 transfer results are reported in their dedicated
sections so that their different detectors, splits, and objectives are not
pooled with this modern evaluation.
